\section{Conclusion}
\label{sec:conclusion}
This paper presented an open and fully reproducible benchmark for day-ahead hourly load forecasting of the four subsystems of the Brazilian interconnected system, built exclusively from public operator data and evaluated over almost three years with an operationally realistic information set. Linear ARX models, a gradient boosting model, a neural network, and their average were compared, and the best open models were evaluated against the official day-ahead forecast of the system operator.

The main findings are as follows. First, a simple average of ARX, GBM, and MLP forecasts, without any weather input, was 16\% more accurate than the official forecast on average (MAPE 2.44\% versus 2.90\%), with significant gains in three of the four subsystems. Second, the official forecast shows a systematic midday under-forecast that grew from 2024 to 2026 together with the share of behind-the-meter generation, whereas the open models remained nearly unbiased; this bias is the main source of the difference in accuracy, although the open models remained more accurate after the bias was reduced in mid-2026. Third, under drift, the calibration window length matters more than the refitting frequency, and level normalisation is essential for tree-based models. Fourth, adaptive conformal inference with delayed feedback produced calibrated prediction intervals, whereas quantile gradient boosting was strongly overconfident.

Future work will add operational weather forecasts, extend the comparison to other operators that publish their official forecasts, and study forecast combinations that explicitly correct the hour-specific bias of operational forecasts. The code and data pipeline are released so that the benchmark can be updated as new data are published.
